

\documentclass[11pt]{article}
\usepackage[utf8]{inputenc}
\usepackage{amsmath}
\usepackage{lipsum} % to generate some filler text
\usepackage{fullpage}
\usepackage{mathrsfs}
\usepackage{svg}
\usepackage{xr}
\usepackage{amsthm,cite,url,graphicx,booktabs,lipsum,color,bm,caption,subcaption,soul}
\usepackage{pifont,tikz,paralist,multirow,amssymb}
\usepackage{xifthen}
\usepackage{enumerate}
\usepackage{titlesec}
\newcounter{reviewer}
\setcounter{reviewer}{0}
\newcounter{point}[reviewer]
\setcounter{point}{0}

% This refines the format of how the reviewer/point reference will appear.
%\renewcommand{\thepoint}{C\,\thereviewer.\arabic{point}}

\renewcommand{\thepoint}{Comment\,\thereviewer.\arabic{point}}

% command declarations for reviewer points and our responses

\newenvironment{point1}
   {\refstepcounter{point} \bigskip \noindent {\textbf{Associate Editor Comment~}} ---\ }
   {\par }

\newcommand{\reviewersection}{\stepcounter{reviewer} \bigskip \hrule
                  \section*{Reviewer \thereviewer}}

% \newenvironment{point}

%    {\refstepcounter{point} \bigskip \noindent {\textbf{Reviewer~\thepoint} } ---\ }
%    {\par }

\newenvironment{point}
   {\refstepcounter{point} \bigskip \noindent {\textbf{Reviewer~\thepoint} } ---\ }
   {\par }


\newcommand{\shortpoint}[1]{\refstepcounter{point}   \bigskip \noindent 
	{\textbf{Reviewer~Point~\thepoint} } ---~#1  }

\newenvironment{reply}
   {\medskip \noindent \color{blue} \begin{sf}\textbf{Reply}:\  }
   {\medskip \par \end{sf}}

\newcommand{\shortreply}[2][]{\medskip \noindent \begin{sf}\textbf{Reply}:\  #2
	\ifthenelse{\equal{#1}{}}{}{ \hfill \footnotesize (#1)}%
	\medskip \end{sf}}

\makeatletter
\newcommand*{\addFileDependency}[1]{% argument=file name and extension
  \typeout{(#1)}
  \@addtofilelist{#1}
  \IfFileExists{#1}{}{\typeout{No file #1.}}
}
\makeatother

\newcommand*{\myexternaldocument}[1]{%
    \externaldocument{#1}%
    \addFileDependency{#1.tex}%
    \addFileDependency{#1.aux}%
}

\myexternaldocument{main}
%\externaldocument[I-]{main}


\def\thesubsubsection{\arabic{subsubsection}}
\titleformat{\subsubsection}[runin]
{\normalfont\bfseries}
{\indent\thesubsubsection)}{0.5em}{}


\begin{document}
\section*{Notes on revision made to manuscript: Fairness and Transparency in ML: A Methodological Framework for Ethical Model Evaluation
}
\vspace{10pt}

The authors would like to thank the editor and reviewers for their constructive comments and suggestions that have helped improve the quality of this manuscript.
% Let's start point-by-point with Reviewer 1

\section*{Response to the reviewers}


%%%%%%%%%%%%%%%%%%%%%%%%%%%%%%%%%%%%%%%%%%%%%%%%%%%
\reviewersection
% General intro text goes here

% Point one description 

%%%%%%%%%%%%%%%%
\begin{point}
	The structure of Section 7 is a bit confusing: there are too many figures and tables that might be better placed into an Appendix, as it's hard to follow the flow of the paper with all the breaks in the text. It would be better to aggregate the results of the improved and base models in a more synthetic way to better highlight the key takeaways. As currently presented, assessing the benefit brought by the proposed methodology is really not straightforward. \label{comment_11}
\end{point}

% Our reply
\begin{reply}

Thank you for your detailed feedback. To address your concerns, we have restructured the manuscript as follows:

    \begin{enumerate}
        \item The original Section 3 has been renamed from ``Experiments \& Results” to ``Results \& Analysis” to better align with its content and focus.
        \item The subsection on experiments (previously Section 3.1) has been relocated to the Methodology section (now as Subsection 2.4). This adjustment aims to provide a more cohesive presentation of the experimental setup alongside the methodology.
        \item The original Section 4 (Discussion) has been removed. Instead, we have integrated the explanation and discussion of results directly into the corresponding subsections (3.1 and 3.2) of the revised ``Results \& Analysis” section, to ensure a more streamlined narrative.
    \end{enumerate}

We believe these modifications should significantly enhance the clarity and readability of the manuscript, as well as address your concerns regarding the structure and synthesis of the results. 
	   
\end{reply}

\begin{point}
	I'm not entirely convinced about the choice of the age distribution (a variable that was not available in the original data and that the authors synthetically added). The choice is based on a very coarse information about the crime age prevalence; I wonder why not expand the analysis using different distributions as well - at least to empirically demonstrate that this synthetic data is not skewing the results. \label{comment_12}
\end{point}

% Our reply
\begin{reply}
	Thank you for raising this point. To address your concern, we performed an additional simulation using a uniform distribution for the age variable, in contrast to the original normal distribution. The results of this experiment were practically identical to those obtained with the normal distribution, demonstrating that the choice of distribution does not skew the results. To maintain the flow and clarity of the main text, the results of this additional experiment have been included in Appendices A and B. We hope this supplementary analysis provides sufficient evidence to support the robustness of our experiments.
\end{reply}

\begin{point}
    The methodology is claimed by not fully explained nor justified. This is noted first by observing that there is not methodology section (not even a brief one) in the paper. The actual methodology is presented in Figure 1 with scant details provided in the text. In the current state of the paper is not really possible to understand if the methodology is technically sound. It is also not clear whether it can be applied to other use cases.
\end{point}

\begin{reply}
    Thank you for your observation. The methodology has been detailed in Section 2 of the manuscript, which existed in the previous version as well. However, we believe the changes made to the structure of the manuscript, as described in the response to Comment 1.1, improve the overall clarity and flow, making the methodology easier to understand. Specifically, moving the experimental details to the Methodology section and integrating discussions into the Results \& Analysis section provide a more cohesive narrative.
    
%    We kindly invite you to review the updated manuscript and let us know if further adjustments are needed to clarify the methodology.
\end{reply}

\begin{point}
    What is intended with "level of development of the crime"? See Sec. 2.2.
\end{point}

\begin{reply}
    Thank you for pointing this out. The term “level of development of the crime” referred to the stage or degree of development in which the crime was committed, such as whether it was consummated, attempted, or frustrated. We acknowledge that “level of development” was a poor translation, and we have since replaced it with “crime stage,” which more accurately conveys the intended meaning. This change has been implemented throughout the manuscript for clarity.
\end{reply}

\begin{point}
    What do you mean with "set" in the following sentence "A set can have one or several outcomes (i.e. “Outcome 1”, “Outcome 2”, etc)"? It's not clear what different outcomes correspond to -- are these the different outcomes provided at different stages of the process? This should be made clearer.
\end{point}

\begin{reply}
     Thank you for your observation. The term “set” was unclear and has been corrected to “dataset” for clarity. Regarding the outcomes, we revised the paragraph to provide a more detailed explanation. Briefly, outcomes represent the results of legal proceedings, which may change over time through different stages of the judicial process (e.g., from a prison sentence to house arrest after an appeal). To simplify the analysis, only the most recent outcome was considered: the second outcome for the drug trafficking dataset and the first (and only) outcome for the petty theft dataset. We hope this clarification resolves your concerns. 
\end{reply}


\begin{point}
    Can the proposed methodology be applied more broadly than the use cases considered in the paper?
\end{point}

\begin{reply}
    Thank you for your question. Yes, the proposed methodology is agnostic to the specific problem and the type of model used. It is designed to be applicable to any classification problem where ethical considerations are necessary. The framework can be adapted to different domains and datasets, making it broadly usable beyond the use cases considered in the paper. This clarification has been added to the conclusion to ensure it is clear and accessible for the reader.
\end{reply}


\begin{point}
    How hyperparameter tuning is performed?
\end{point}

\begin{reply}
    Thank you for your question. The hyperparameter tuning for the models was conducted as follows: For the MLP model, we used manual search to identify the optimal configuration. For the LightGBM model, we used a combination of manual search and grid search was used to refine the hyperparameters. This clarification has been added to both experiments.
    \\
    \\
    For the improved versions of both models, some additional hyperparameter tuning was performed, while also applying the weights suggested by the reweighting function from AIF360 and the compute\_class\_weight function from Scikit-learn to address class and sample imbalances.
\end{reply}


\begin{point}
    Is there any theoretical justification for the methodology steps? What are the guiding principles that demonstrate that the methodology is valid?

\end{point}

\begin{reply}
    Although the proposed methodology does not have a theoretical justification, it does follow a typical data science life cycle. The main difference is that we incorporate pre- and post-training stages that measure levels of bias and disparities, and therefore the model selection process is not guided in the first instance by the usual metrics such as accuracy or f1 score, but by the metrics associated with bias and disparities. 
\end{reply}


\begin{point}
    In view of the lack of justification, it is really hard to understand why some choices were made; I'm not claiming that these choices are wrong, but I'm trying to understand their motivation. E.g., Why in 3.2 a multi-layer perceptron was used?
\end{point}

\begin{reply}
    Thank you for your observation. It is important to note that the specific choice of models is not central to this research, as the focus lies on the proposed framework rather than the model architectures themselves. Therefore, there was no theoretical justification for selecting this specific architecture beyond its practical suitability for the task at hand. 
    
    In particular, the choice of using a multi-layer perceptron (MLP) in Section 3.2 was primarily driven by its adaptability to the specific problem provided by Chile's Public Criminal Defense Office (DPP), which is to classify whether a defendant would have a favourable or unfavourable outcome to its trial.
    
    While the DPP required models to address their operational needs, the analysis and development conducted were independently led as part of this research, and are not being used in production by them.
\end{reply}


\begin{point}
    Why weren't other ML models tried as well?\end{point}

\begin{reply}
    Thank you for your comment. The focus of this research is not on the specific machine learning models themselves, but rather on the application and evaluation of the proposed framework. 
    \\
    \\
    While other types and architectures of classification models were evaluated, we ultimately selected the two used in the study due to their ease of development, facilitated by the availability of robust frameworks such as TensorFlow and Scikit-learn. Future work could involve testing additional models to further validate the framework’s flexibility and robustness.
\end{reply}

\begin{point}
    Why two different MLP's configurations were used in the two cases (the base one and the improved one) described in Sec. 3.2? Shouldn't the 2 MLPs have the same configuration to ensure a fair comparison?

\end{point}

\begin{reply}
    Thank you for your observation. While it is true that different configurations of the MLP were used in the two experiments, the changes did not result in significant differences in the model’s performance. As mentioned in the text, the adjustments to the configuration were made to conduct additional tests and explore the effects of varying parameters, and is part of the model improvement process. This is also true for Experiment 3.2 (petty theft), where the configuration of the improved model — built with LightGBM — was also modified as part of this iterative improvement process. 
\\
\\
    We acknowledge that keeping the configurations identical could enhance the fairness of the comparison, and this is a valuable point for future considerations. However, based on our results, we can confirm that the conclusions drawn from the experiments remain valid regardless of these adjustments. Nevertheless, a clarification has been added to that part of the experiment.
\end{reply}

On the minor remarks and typos:

\begin{point}
    Considering the focus on the AI lifecycle, the following papers could be useful related works:
- Calegari R, Castañé GG, Milano M, O'Sullivan B. Assessing and enforcing
fairness in the AI lifecycle. IJCAI International Joint Conference on
Artificial Intelligence.
- Agarwal A, Agarwal H. A Seven-Layer Model for Standardising AI Fairness Assessment. arXiv preprint arXiv:2212.11207. 2022 Dec 21.
- Kasy M, Abebe R. Fairness, equality, and power in algorithmic decision-making. InPro

\end{point}

\begin{reply}
    Thank you for the suggestion. The mentioned papers have been added to the manuscript, and can be found in section 1.1 (Background).
\end{reply}

\begin{point}
    Figures 11, 15, 16, 17 contain non-English text; I think cannot be accepted in an international journal. The authors should re-create the picture translating the non-English words. I saw that a translation was provided in Sec. 2.2 but I'm still skeptical.

\end{point}

\begin{reply}
    Thank you for pointing this out. This has been corrected. Every figure in the manuscript should now have its axis labels and feature names in English.
\end{reply}

\begin{point}
    "How can we ensure that these models do not possess risks to our society?" --> maybe a better sentence would have been "How can we ensure that these models do not pose risks to our society"

\end{point}

\begin{reply}
    Thank you for the suggestion. The sentence has been revised to improve clarity as proposed.
\end{reply}

\begin{point}
   "This is, models cannot “reason” in the same manner as a human" --> "That is, models cannot “reason” in the same manner as a human"

\end{point}

\begin{reply}
    Thank you for pointing this out. The sentence has been updated accordingly.
\end{reply}

\begin{point}
    "This is largely due to its recent emergence and a lack of obligation for their use by private companies" --> "This is largely due to their recent emergence and a lack of obligation for their use by private companies"

\end{point}

\begin{reply}
    Thank you for your comment. This paragraph in particular has been changed to include the papers suggested in Comment 1.12.
\end{reply}

\begin{point}
    ``These contains the records of multiple individuals accused" --> "These contain the records of multiple individuals accused"

\end{point}

\begin{reply}
    Thank you for identifying this issue. The suggested correction has been implemented in the manuscript.
\end{reply}

\begin{point}
    ``For the second dataset, we use the data from the Centre for Crime Studies and Analysis (de Estudios y Análisis del Delito, ) for the crime of theft between the years 2017 and 2022. This variable also considers a normal distribution"

\end{point}

\begin{reply}
    Thank you so much for pointing this out. We have corrected that format error, as well as some others found in section 1.1 (Background).
\end{reply}

\begin{point}
    I'm not sure that scikit-learn can be placed under the umbrella of "optimisation tools", it's rather the go-to tool for ML. Similarly, AI Fairness 360 and Fairlearn would better be grouped together with Aequitas under a category such as "Fairness Enhancers"

\end{point}

\begin{reply}
    Thanks for pointing this out. Scikit-Learn has been moved to the first category (essential development tools); while AIF360 and Fairlearn were moved to ``Fairness Enhancers".
\end{reply}

%%%%%%%%%%%%%%%%%%%%%%%%%%%%%%%%%%%%%%%%%%%%%%%%%%%
\reviewersection

%%%%%%%%%%%%%%%%
\begin{point}
\label{comment_21} As stated by the authors, due to the lack of many critical variables (such as the age, sex, and previous cases of the defendant, among others), they applied a normal distribution to create a newly added column “age” to the first dataset (drug trafficking) based on other data about drug trafficking. In a similar way, they added a column of age to the second dataset (petty theft). The preprocessing method is reasonable; however, there are issues causing the mismatching of the generated column with the real situation, though it is impossible to know the truth in this study. Generative models have been used for more accurate synthetic data generation. Alternatively, if there could be overlapping columns of both datasets, a synthetic column of age could be generated based on correlation with other columns; in this way, the correlation with other attributes might be closer to reality.
\end{point}

\begin{reply}
 Thank you for your thoughtful and constructive feedback. We understand your concerns regarding the use of a synthetic age column and its potential mismatch with the real situation. To address this, we conducted an additional analysis where we used a uniform distribution to generate the synthetic age column as an alternative to the normal distribution. The results from this analysis were consistent with those obtained using the normal distribution, suggesting that the synthetic age variable does not skew the results significantly.
\\
\\
We acknowledge that generative models could provide a more robust approach to synthetic data generation, particularly by capturing correlations between attributes. However, due to the limitations of the datasets and the scope of this study, we focused on simpler methods for creating synthetic columns. We appreciate your suggestion regarding using correlations with other columns, and we will consider this approach for future work to enhance the realism of the synthetic data.
\\
\\
We have included the results of the additional analysis in Appendices A and B for reference. 
\end{reply}


%%%%%%%%%%%%%%%%
\begin{point}
\label{comment_22} The proposed framework in this manuscript integrates multiple tools and techniques to achieve the goal, such as SHAP for XAI, sampling, normal distribution for data preprocessing, fairness disparity as evaluation metrics, and using accuracy, precision, recall, and F1-score for classification evaluation. AIF360’s reweighting function and Scikit-Learn’s compute class weight method have also been used regarding re-weighting and balance. Based on integrated results, extensive experimental results have been presented in this manuscript. However, the results and analysis themselves are reasonably stated, but not well presented and coordinated to accompany each other. In the earlier part of the experiments, the authors presented an extensive amount of experimental results in tables and figures. All the discussion and analysis of the results come intensively later. It is strongly suggested that the key analysis and discussion on specific tables and figures be put around the visualized resulted to facilitate reading. In the later part, some summaries and highlights based on the discussion could be added if need.
\end{point}

\begin{reply}
 Thank you so much for your feedback. To address your concerns, we have restructured the manuscript to improve the presentation and coordination of results and analysis. Specifically:

 \begin{enumerate}
        \item The original Section 3 has been renamed to ``Results \& Analysis” to reflect the integration of results and their corresponding discussion, facilitating a more cohesive narrative.
        
        \item The subsection on experiments (previously Section 3.1) has been relocated to the Methodology section (now as Subsection 2.4). This adjustment aims to provide a more cohesive presentation of the experimental setup alongside the methodology.
        \item The original Section 4 (Discussion) has been removed, and the discussion and analysis of results have been incorporated directly into the subsections of the “Results \& Analysis” section. This adjustment ensures that key observations and highlights are provided alongside the visualized results, improving readability and flow.
    \end{enumerate}

We believe these changes address your suggestion to streamline the presentation of results and discussion. 

\end{reply}
%%%%%%%%%%%%%%%%
\begin{point}
The authors stated that the proposed framework evaluates the base model regarding its performance, biases (through counterfactual analysis), and fairness (disparities). Fairness (disparities) evaluations have been conducted through the experiments. However, there are no explicit explanations about how the mentioned counterfactual analysis has been conducted.
\end{point}

\begin{reply}
 Thank you for your observation. The mention of counterfactual analysis in that sentence has been removed, as we were ultimately unable to conduct the counterfactual analysis due to the lack of necessary variables and data required for this type of evaluation. However, this remains an important direction for future work, as it could complement the other bias metrics analyzed in the experiments.
\end{reply}
%%%%%%%%%%%%%%%%
\begin{point}
“Region” as a protected attribute has been tested in the extensive experiments. More protected attributes are expected to be tested to validate the proposed framework. Instead of age, gender, etc., some other attributes that are not very sensitive but still can provide grouping of the defendants could be considered in the future.\label{comment_31}
\end{point}

\begin{reply}
 Thank you for your constructive feedback. We completely agree that testing additional protected attributes would further validate the proposed framework. However, due to the limitations of the available data, “Region” was the only variable that was both viable and meaningful for conducting the analysis in this study. Nevertheless, this point has been added to the future work section in the conclusions.
\end{reply}

\begin{point}
Through the whole experimental results, filled with the SHAP waterfall diagrams (e.g., Figures 2, 3) and the numeric results (e.g., Tables 2, 3), the SHAP value distribution could be more informative. For example, following each attribute name on the y-axis, an average value could be added. For each attribute, the corresponding value distribution, positive and negative values, could be marked on it. Although the values are also in the table, the delivered information from figure and table could be more balanced. The waterfall diagrams are very visible but not informative; the values in the table are detailed but hard to read.\label{comment_31}
\end{point}

\begin{reply}
 Thank you for your constructive feedback. In response to your suggestions, we have enhanced the SHAP beeswarm plot by including the average SHAP value for each attribute directly in the y-axis labels. This addition provides immediate context for the contribution of each feature, making the visualization more informative and accessible.
\\
\\
Additionally, we have consolidated the two previous tables into a single table that presents all relevant metrics for both the base and improved models. This unified table includes the average SHAP values and average absolute SHAP values. We hope that by merging these metrics into one table, we provide a more comprehensive and streamlined summary of the SHAP results, ensuring clarity and ease of comparison. We hope these updates address your concerns and improve the overall presentation.
\end{reply}

\begin{point}
In Section 1.2, there are explanations of the terminologies. For disparity, there should be more information, especially with examples, since the relevant terms and metrics are used throughout the experiments. More basic and brief information about FPR, FNR, FDR, and FOR along with the confusion matrix is also useful. Similarly, there is a need for more basic explanations of the three groups including Reference Group, Major Group, and Min. Metrics Group. This will pave the foundation to facilitate reading in the experimental analysis.\label{comment_31}
\end{point}

\begin{reply}
 Thank you for your thoughtful comment. In response, we have enhanced the results section by adding a detailed introductory explanation just before the analysis begins. This new section includes the definitions for FPR, FNR, FDR, and FOR, alongside a brief explanation of the roles of the Reference Group, Major Group, and Min. Metrics Group, which now should be clearly defined to ensure a solid understanding of the terminology used throughout the experiments.
\end{reply}

\begin{point}
For those figures (e.g., 12-14) which present model assessment for the different groups, more explicit explanations are required. Though clearly the three colors denote the parity test (reference, pass, fail), it is unclear about the essential meaning of the different colors and how information is embodied via the color with different lengths of the line.\label{comment_31}
\end{point}

\begin{reply}
 Thank you so much for this comment. To address your concern, we have added detailed explanations of the figures to subsection 3.1.4. These explanations should clarify the meaning of the colours and bar lengths used in the visualisations. 
\end{reply}

\begin{point}
In some tables, English attributes are presented, but most of the figures and tables are filled with Spanish. It is recommended that the authors translate them into English for consistency.\label{comment_31}
\end{point}

\begin{reply}
 Thank you for highlighting this inconsistency. All tables and figures have been reviewed, and the Spanish text has been translated into English to ensure consistency throughout the manuscript. 
\end{reply}

\begin{point}
A few typos, such as a redundant “the” in “as the it” on page 49, and double quotation mismatching for “Region” on page 44.\label{comment_31}
\end{point}

\begin{reply}
    Thank you for pointing this out. We have corrected these typos and other format errors found through the manuscript.
\end{reply}

\begin{point}
The fonts within several figures are too blurry to read (e.g., Figures 4-6).\label{comment_31}
\end{point}

\begin{reply}
 Thank you for your comment. Regarding the specific cases mentioned (Figures 4–6) and all other figures generated by Aequitas, we regret to inform that it is not possible to adjust the font size or color as these are directly generated by Aequitas’ built-in functions. Unfortunately, the library does not provide customisation options for these graphical elements.
\end{reply}


%\bibliographystyle{IEEEtran}
%\bibliography{IEEEabrv,mybibfile}
\end{document}


\begin{point}

\end{point}

\begin{reply}

\end{reply}

